\documentclass[11pt]{article}
\usepackage[T1]{fontenc}
\usepackage{lmodern}
\usepackage{microtype}
\usepackage[margin=1in]{geometry}
\usepackage{amsmath,amssymb,mathtools,amsthm}
\usepackage{graphicx}
\graphicspath{{figures/}{/mnt/data/}}
\usepackage{booktabs}
\usepackage{tabularx}
\usepackage{array}
\usepackage{enumitem}
\usepackage{xcolor}
\usepackage{hyperref}
\usepackage[nameinlink,noabbrev]{cleveref}
\usepackage{caption}
\usepackage{float}
\usepackage{setspace}
\usepackage{fancyhdr}
\usepackage{titlesec}

\definecolor{linkblue}{RGB}{24,74,128}
\hypersetup{colorlinks=true,linkcolor=linkblue,citecolor=linkblue,urlcolor=linkblue,pdftitle={Environment-Induced Quantum Learning - Version 3}}
\setstretch{1.04}
\setlength{\parindent}{0pt}
\setlength{\parskip}{0.48em}
\setlist[itemize]{topsep=0.25em,itemsep=0.15em,leftmargin=1.6em}
\setlist[enumerate]{topsep=0.25em,itemsep=0.15em,leftmargin=1.8em}
\titleformat{\section}{\large\bfseries}{\thesection}{0.55em}{}
\titleformat{\subsection}{\normalsize\bfseries}{\thesubsection}{0.5em}{}

\setlength{\headheight}{14pt}
\pagestyle{fancy}
\fancyhf{}
\fancyhead[L]{Environment-Induced Quantum Learning}
\fancyhead[R]{Version 3 - theory + simulation}
\fancyfoot[C]{\thepage}
\renewcommand{\headrulewidth}{0.3pt}

\newtheorem{theorem}{Theorem}
\newtheorem{proposition}{Proposition}
\newtheorem{lemma}{Lemma}
\newtheorem{corollary}{Corollary}
\newtheorem{definition}{Definition}
\newtheorem{remark}{Remark}

\newcommand{\tr}{\operatorname{Tr}}
\newcommand{\supp}{\operatorname{supp}}
\newcommand{\TV}{\operatorname{TV}}
\newcommand{\cL}{c_L}
\newcommand{\eps}{\varepsilon}
\newcommand{\Mcal}{\mathcal{M}}
\newcommand{\Zcal}{\mathcal{Z}}
\newcommand{\Xcal}{\mathcal{X}}
\newcommand{\Oeps}{\mathcal{O}_{\eps,L}^{\mathfrak M}}

\begin{document}

\begin{center}
{\LARGE\bfseries Environment-Induced Quantum Learning (EIQL)}\\[0.45em]
{\large Self-Supervised Measurement Discovery from Redundant Quantum Records}\\[0.7em]
{\normalsize Version 3 - Theory and Simulation Manuscript}\\
{\small 14 August 2026}
\end{center}

\vspace{0.5em}
\begin{abstract}
Environment-Induced Quantum Learning (EIQL) is a quantum measurement-learning framework in which redundancy across independently accessible environmental fragments supplies the self-supervised signal for discovering an initially unknown local decoder. The abstract principle of agreement across multiple views is not claimed to be new; the quantum operational distinction is that the representation is not an already available classical feature map but must be physically instantiated by choosing a POVM on each fragment. EIQL therefore asks which hardware-constrained local measurements produce outputs that are simultaneously rich and mutually consistent, without target labels, without a specified target measurement, and without requiring direct access to the system.

The theory identifies pointer information under learned local decoders rather than a unique physical POVM. Under exact spectrum broadcast structure (SBS), exact rich agreement forces every local output to be a deterministic coarse-graining of the pointer variable; when the output alphabet and hardware class resolve the record supports, the pointer information is recovered up to relabeling. For states close to SBS, high richness and low disagreement imply a quantitative residual-uncertainty bound and a simple explicit Bayes-recovery certificate. A finite-shot proposition certifies pre-specified finite decoder classes. Theory-matched five-qubit simulations reproduce an analytic independent-product null ceiling and remain separated from the null under a simple NISQ-style noise model. A second benchmark, based on the published six-photon Quantum-Darwinism architecture of Chen et al., hides an independent local basis on every environmental record: EIQL recovers the strong-record decoders, identifies the three reliable fragments in the mixed-quality setting, and approaches the physical disagreement floor. Finally, an environment-only pair-moment implementation is compared at equal copy budget with a stronger system-assisted correlation baseline and with full Pauli tomography in measurement-setting count. These results support EIQL as a specific quantum learning formulation; they do not establish quantum advantage or universal resource superiority.
\end{abstract}

\textbf{Status and scope.} This Version 3 manuscript freezes the first-paper scope at theory plus simulation. The mathematical statements are conditional theorems under explicit assumptions, and all numerical results are simulation results unless stated otherwise. No quantum advantage, hardware advantage, or unique-POVM recovery is claimed. A real-hardware EIQL demonstration is left as future work rather than treated as a prerequisite for the present paper. Literature positioning has been checked through 14 August 2026, with particular attention to Quantum Darwinism experiments, unknown-measurement learning, self-supervised quantum learning, and unsupervised observable discovery.

\tableofcontents
\newpage

\section{Research question and scope}

Quantum Darwinism treats the environment as a communication channel that proliferates records of selected system observables. In the environment-as-witness formulation, redundant imprints are associated with preferred pointer observables selected by decoherence and einselection \cite{ollivier2004,ollivier2005,zurek2003}. Spectrum broadcast structure (SBS) gives a particularly sharp state-level form of objective records: a classical system variable is correlated with locally distinguishable states in multiple subenvironments \cite{horodecki2015,le2019}.

EIQL starts one step later. It assumes that an observer has repeated physical access to environmental fragments, but is \emph{not} told which fragment measurement reveals the objective record. The central learning question is:

\begin{quote}
Given repeated preparations of quantum environmental fragments and a physically implementable measurement class, can an observer learn a local measurement/decoder whose outputs are both informative and redundantly consistent across fragments, without labels and without prior knowledge of the pointer observable?
\end{quote}

This question is deliberately narrower than ``discovering objectivity'' in general. Redundancy, pointer selection, uniqueness of redundant imprints, and observer consensus already have substantial prior art \cite{ollivier2004,fu2021,touil2025}. EIQL instead treats the \emph{decoder itself} as the learned object.

The learning principle also has classical relatives: multiview agreement, common-information extraction, and self-supervision all predate EIQL, and quantum self-supervised learning exists as a broader research direction \cite{gacs1973,jaderberg2021}. The intended claim is therefore not a new universal learning paradigm. The operational quantum difference is that the views are quantum states and the learner must choose the physical measurement that creates the classical representation. Incompatible candidate measurements cannot all be read simultaneously from the same copy. This is why the present paper is framed as a \emph{quantum measurement-learning framework/problem formulation}, not merely as a quantum implementation of a classical agreement loss and not as a claim of generic quantum learning advantage \cite{ciliberto2018}.

The framework is also distinct from the usual variational-QML pattern
\[
\rho \longrightarrow U(\theta) \longrightarrow \text{measurement} \longrightarrow \text{loss} \longrightarrow \text{optimize }\theta.
\]
A basic EIQL implementation can search a shallow local measurement family directly, so deep parameterized circuits and their gradient landscapes are not necessary. This does not constitute a theorem that EIQL can never exhibit optimization pathologies; it only means that barren plateaus are not intrinsic to the minimal formulation. Barren plateaus remain an established concern for broad families of variational quantum circuits \cite{mcclean2018,larocca2024}.

\section{Physical setting}

Let $S$ be a finite-dimensional system and let the environment be partitioned into $m\ge 2$ fragments,
\[
E = E_1\otimes E_2\otimes\cdots\otimes E_m.
\]
The learner receives repeated preparations of a state $\rho_{SE_1\cdots E_m}$ but need not access $S$ directly. On each fragment $E_j$, it may apply a POVM
\[
M_j=\{M_{j,z}\}_{z\in\Zcal},\qquad M_{j,z}\succeq 0,\qquad \sum_{z\in\Zcal}M_{j,z}=I_{E_j},
\]
where $|\Zcal|=L\ge 2$. The resulting classical output is $Z_j\in\Zcal$.

A hardware restriction is represented by a set $\mathfrak M_j^{(L)}$ of allowed $L$-outcome POVMs. A local \emph{decoder} is a pair $\mathcal D_j=(M_j,\pi_j)$, where $M_j\in\mathfrak M_j^{(L)}$ and $\pi_j$ is an arbitrary classical permutation of the output alphabet $\Zcal$. Output relabelings are free: the operational variable is
\[
Z_j=\pi_j(Y_j),
\]
where $Y_j$ is the raw detector label. The joint admissible decoder class contains all such relabelings of
\[
\mathfrak M^{(L)}=\mathfrak M_1^{(L)}\times\cdots\times\mathfrak M_m^{(L)}.
\]
This convention prevents a purely notational mismatch---for example, one fragment outputting $X$ and another outputting $1-X$---from being counted as physical disagreement. The hardware class can encode single-qubit projective measurements, a finite set of calibrated bases, bounded-depth local pre-rotations followed by computational-basis measurement, or a restricted family of local instruments.

\begin{definition}[Operational decoder equivalence]
Fix a relevant record family $\mathcal R_j=\{\sigma_x^{E_j}:x\in\Xcal\}$. Two local decoders $\mathcal D_j$ and $\mathcal D'_j$ are operationally equivalent on $\mathcal R_j$, written $\mathcal D_j\sim_{\mathcal R_j}\mathcal D'_j$, if they induce the same conditional output statistics on every record state:
\[
P_{\mathcal D_j}(Z=z\mid\sigma_x^{E_j})=P_{\mathcal D'_j}(Z=z\mid\sigma_x^{E_j})
\qquad\forall x,z.
\]
EIQL aims to identify pointer information and decoder behavior on the physically occupied record states. It does not, without additional assumptions, identify a unique POVM operator on the full Hilbert space.
\end{definition}

\subsection{Spectrum broadcast structure}

The reference ideal is a perfect SBS state
\begin{equation}
\sigma_{SE_1\cdots E_m}
=\sum_{x\in\Xcal}p_x\,|x\rangle\!\langle x|_S\otimes
\sigma_x^{E_1}\otimes\cdots\otimes\sigma_x^{E_m},
\label{eq:sbs}
\end{equation}
with $p_x>0$ on the support and local distinguishability
\begin{equation}
\sigma_x^{E_j}\sigma_{x'}^{E_j}=0\qquad (x\ne x')
\label{eq:orth}
\end{equation}
for every fragment $j$. The classical random variable $X$ with $P(X=x)=p_x$ is the broadcast pointer variable. This is the classical information made objective by the SBS structure; the records that carry it are quantum states \cite{horodecki2015,korbicz2021}.

For any tuple of local measurements, the product form of \cref{eq:sbs} gives conditional independence
\begin{equation}
P_\sigma(z_1,\ldots,z_m\mid x)=\prod_{j=1}^m P_\sigma(z_j\mid x).
\label{eq:condind}
\end{equation}
This conditional-independence property is the mathematical bridge to common-randomness extraction.

\section{The EIQL objective}

The first draft used average pairwise disagreement and the entropy of a distinguished fragment. That objective can be gamed when $m$ is large: one fragment may carry a rich variable while most other fragments agree on something trivial. Version 2.1 retains the symmetric worst-case quantities introduced in Version 2.

For an admissible measurement tuple $\mathbf M=(M_1,\ldots,M_m)$ define the \emph{worst-pair disagreement}
\begin{equation}
D_\rho(\mathbf M)=\max_{1\le i<j\le m}P_\rho(Z_i\ne Z_j),
\label{eq:disagree}
\end{equation}
and the \emph{redundant richness}
\begin{equation}
R_\rho(\mathbf M)=\min_{1\le j\le m}H_\rho(Z_j).
\label{eq:richness}
\end{equation}
The minimum in \cref{eq:richness} prevents a single high-entropy fragment from defining the representation.

\begin{definition}[Hardware-constrained EIQL value]
For tolerance $\eps\in[0,1]$, alphabet size $L$, and hardware class $\mathfrak M^{(L)}$, define
\begin{equation}
\mathcal O_{\eps,L}^{\mathfrak M}(\rho)
=
\sup_{\mathbf M\in\mathfrak M^{(L)}}R_\rho(\mathbf M)
\quad\text{subject to}\quad
D_\rho(\mathbf M)\le \eps.
\label{eq:objective}
\end{equation}
If the feasible set is empty, $\mathcal O_{\eps,L}^{\mathfrak M}(\rho)$ is undefined (or may be assigned $-\infty$ by convention).
\end{definition}

The quantity $\eps\mapsto\mathcal O_{\eps,L}^{\mathfrak M}(\rho)$ is the \emph{agreement-richness frontier}. It describes how much common information is operationally extractable at a prescribed worst-fragment reliability.

For selecting an operational decoder representative, \cref{eq:objective} may be lexicographically refined: first maximize $R_\rho$, then among solutions within a small tolerance $\tau$ of the optimum, minimize $D_\rho$. This tie-break does not change the frontier itself; it selects a stable representative when richness is flat over many measurements.

\section{Exact pointer-information identifiability under SBS}

The exact result is best interpreted as a foundational lemma, not the primary novelty claim. Related uniqueness and redundant-imprint results are already known in Quantum Darwinism \cite{ollivier2004,fu2021}.

\begin{theorem}[Exact SBS pointer-information structure]
\label{thm:exact}
Let $\sigma$ be a perfect SBS state of the form \cref{eq:sbs}, and let local measurements produce $Z_1,\ldots,Z_m\in\Zcal$. If
\begin{equation}
P_\sigma(Z_1=Z_2=\cdots=Z_m)=1,
\label{eq:exactcons}
\end{equation}
then there exists a deterministic map $g:\Xcal\to\Zcal$ such that
\begin{equation}
Z_j=g(X)\qquad\text{almost surely for every }j.
\label{eq:gX}
\end{equation}
Consequently,
\begin{equation}
H(Z_j)=H(g(X))\le \max_{f:\Xcal\to\Zcal}H(f(X))\le \min\{H(X),\log_2L\}.
\label{eq:alphabetupper}
\end{equation}

If, in addition, each hardware class $\mathfrak M_j^{(L)}$ contains an SBS-support-resolving measurement followed by arbitrary classical post-processing into $\Zcal$, then
\begin{equation}
\mathcal O_{0,L}^{\mathfrak M}(\sigma)
=
\max_{f:\Xcal\to\Zcal}H(f(X)).
\label{eq:exactopt}
\end{equation}
In particular, if $L\ge |\supp X|$, then
\begin{equation}
\mathcal O_{0,L}^{\mathfrak M}(\sigma)=H(X),
\end{equation}
and every exact-consensus maximizer with entropy $H(X)$ recovers $X$ up to a one-to-one relabeling on its support.
\end{theorem}

\begin{proof}
Because the average of the conditional probabilities of the event in \cref{eq:exactcons} equals one, that event has conditional probability one for every $x$ with $p_x>0$. Consider fragments $1$ and $2$. Write
\[
p_x(z)=P_\sigma(Z_1=z\mid X=x),\qquad q_x(z)=P_\sigma(Z_2=z\mid X=x).
\]
By \cref{eq:condind},
\[
1=P_\sigma(Z_1=Z_2\mid X=x)=\sum_z p_x(z)q_x(z)
\le \max_z p_x(z)\sum_zq_x(z)\le1.
\]
Therefore $\max_zp_x(z)=1$, so $Z_1$ is deterministic given $X=x$. The same argument applies to $Z_2$, and exact equality forces their deterministic values to coincide. Repeating with the remaining fragments yields a single function $g$ such that \cref{eq:gX} holds.

The entropy bounds follow because deterministic processing cannot increase Shannon entropy and because an $L$-symbol variable has entropy at most $\log_2L$.

Under the stated hardware expressivity condition, each fragment can first distinguish the orthogonal supports in \cref{eq:orth}, recover $x$, and then output any chosen $f(x)\in\Zcal$. Thus every deterministic coarse-graining $f$ is achievable, proving \cref{eq:exactopt}. If $L\ge|\supp X|$, choose $f$ injective to attain $H(X)$. Conversely, if $H(f(X))=H(X)$, then $H(X\mid f(X))=0$, so $f$ is injective on the support.
\end{proof}

\begin{remark}[Why the alphabet qualification matters]
If $L<|\supp X|$, an $L$-outcome decoder cannot label all pointer values injectively. The exact EIQL optimum is then the maximum-entropy deterministic quantization of $X$, not necessarily $H(X)$. If the allowed hardware class cannot resolve the SBS supports, even that optimum may be unattainable.
\end{remark}

\begin{remark}[Information identifiability is not POVM uniqueness]
The theorem identifies the common classical pointer information encoded in the learned outputs. It does not imply that the physical POVM operator is unique. Distinct POVMs can induce identical statistics on the SBS record supports while differing elsewhere in Hilbert space. The natural learned object is therefore the operational equivalence class $[\mathcal D_j]_{\mathcal R_j}$, together with the induced common variable, rather than a unique operator matrix.
\end{remark}

\section{Robust pointer-information identifiability under approximate SBS}

The central technical result of v2 is the robust theorem. It is stated relative to a nearby perfect SBS state because that supplies a well-defined reference pointer variable $X$.

For an $L$-symbol alphabet define
\begin{equation}
\cL(t)=h_2(t)+t\log_2(L-1),
\label{eq:cL}
\end{equation}
where $h_2$ is binary entropy. On $0\le t\le1-1/L$, $\cL$ is nondecreasing.

\begin{theorem}[Robust pointer-information identifiability under learned local decoders]
\label{thm:robust}
Let $\sigma$ be a perfect SBS state with pointer variable $X$ and let $\rho$ satisfy
\begin{equation}
\frac12\|\rho-\sigma\|_1\le\eta.
\label{eq:traceclose}
\end{equation}
Let $m\ge2$ and let an $L$-outcome measurement tuple $\mathbf M$ satisfy
\begin{align}
D_\rho(\mathbf M)&\le\eps,\label{eq:robdis}\\
R_\rho(\mathbf M)&\ge H(X)-\kappa.\label{eq:robrich}
\end{align}
Define $\delta=\eps+\eta$, and assume $\eta,\delta\le1-1/L$. Then for every fragment $j$,
\begin{equation}
H_\sigma(X\mid Z_j)
\le
\kappa+\cL(\eta)+\cL(\delta).
\label{eq:robustbound}
\end{equation}
Consequently, provided the output alphabet can support this limit, if $\eps,\eta,\kappa\to0$, then the residual uncertainty about the SBS pointer variable from every learned fragment output vanishes.
\end{theorem}

\begin{proof}
Fix a fragment $j$ and select any $k\ne j$. Measurement is a quantum channel, so trace distance contracts. Therefore the classical distributions of $(Z_j,Z_k)$ obtained from $\rho$ and $\sigma$ have total variation distance at most $\eta$. In particular,
\begin{equation}
P_\sigma(Z_j\ne Z_k)\le\eps+\eta=\delta.
\label{eq:sigmadis}
\end{equation}

Conditioned on $X=x$, the SBS product form gives independence of $Z_j$ and $Z_k$. Write
\[
p_x(z)=P_\sigma(Z_j=z\mid X=x),\qquad
q_x(z)=P_\sigma(Z_k=z\mid X=x)
\]
and
\[
\delta_x=P_\sigma(Z_j\ne Z_k\mid X=x).
\]
Then
\[
1-\delta_x=\sum_zp_x(z)q_x(z)\le\max_zp_x(z).
\]
Define
\[
g_j(x)\in\arg\max_zp_x(z),\qquad
e_x=P_\sigma(Z_j\ne g_j(X)\mid X=x).
\]
The preceding inequality implies $e_x\le\delta_x$. Averaging over $X$ gives
\begin{equation}
e:=P_\sigma(Z_j\ne g_j(X))\le\sum_xp_x\delta_x=P_\sigma(Z_j\ne Z_k)\le\delta.
\label{eq:globalerror}
\end{equation}

Let $E=\mathbf 1\{Z_j\ne g_j(X)\}$. Conditional on $(X,E=0)$, $Z_j$ is deterministic; conditional on $(X,E=1)$ it has at most $L-1$ possibilities. Hence
\begin{align}
H_\sigma(Z_j\mid X)
&\le H(E\mid X)+P(E=1)\log_2(L-1)\\
&\le h_2(e)+e\log_2(L-1)\\
&=\cL(e)\le\cL(\delta).
\label{eq:fanoish}
\end{align}
This global error-indicator argument is the key repair relative to v1; it does not require every conditional disagreement $\delta_x$ to be individually small.

Next, the marginal output distributions of $Z_j$ under $\rho$ and $\sigma$ have total variation distance at most $\eta$. The finite-alphabet entropy continuity bound gives
\begin{equation}
|H_\rho(Z_j)-H_\sigma(Z_j)|\le\cL(\eta)
\label{eq:entropycont}
\end{equation}
under the stated range condition \cite{winter2016}. Since \cref{eq:robrich} uses the minimum fragment entropy,
\[
H_\sigma(Z_j)\ge H(X)-\kappa-\cL(\eta).
\]
Finally,
\begin{align}
H_\sigma(X\mid Z_j)
&=H(X)-H_\sigma(Z_j)+H_\sigma(Z_j\mid X)\\
&\le\kappa+\cL(\eta)+\cL(\delta),
\end{align}
which proves \cref{eq:robustbound}.
\end{proof}

\begin{corollary}[Binary output]
For $L=2$,
\begin{equation}
H_\sigma(X\mid Z_j)\le\kappa+h_2(\eta)+h_2(\eps+\eta).
\end{equation}
\end{corollary}

\begin{remark}[Alphabet non-vacuity]
Because every learned output has at most $L$ symbols, $H_\rho(Z_j)\le\log_2L$. Therefore the richness hypothesis $R_\rho(\mathbf M)\ge H(X)-\kappa$ can hold only if
\[
\kappa\ge H(X)-\log_2L
\]
whenever $H(X)>\log_2L$. In particular, the formal limit $\kappa\to0$ requires $H(X)\le\log_2L$. This is a feasibility condition, not an additional proof assumption.
\end{remark}

\begin{corollary}[Operational Bayes recovery]
\label{cor:bayes}
Under the hypotheses of \cref{thm:robust}, define
\[
B=\kappa+\cL(\eta)+\cL(\eps+\eta).
\]
Let $P_{e,j}^*$ be the minimum probability of error for inferring $X$ from the learned classical output $Z_j$. Then
\begin{equation}
P_{e,j}^*\le 1-2^{-B}.
\label{eq:bayesbound}
\end{equation}
Thus $B\to0$ implies an explicit vanishing-error recovery guarantee. The bound is a simple certificate rather than a claim of tight optimality.
\end{corollary}

\begin{proof}
For a fixed observed value $z$, let $p_{\max}(z)=\max_xP(X=x\mid Z_j=z)$ and let $H_z=H(X\mid Z_j=z)$. Shannon entropy dominates min-entropy, so $H_z\ge-\log_2p_{\max}(z)$ and hence
\[
1-p_{\max}(z)\le1-2^{-H_z}.
\]
The function $u\mapsto1-2^{-u}$ is concave. Averaging over $Z_j$ and applying Jensen's inequality gives
\[
P_{e,j}^*=\mathbb E[1-p_{\max}(Z_j)]
\le1-2^{-H(X\mid Z_j)}
\le1-2^{-B}.
\]
\end{proof}

\begin{remark}[Tightness of the Bayes certificate]
The error certificate in \cref{eq:bayesbound} is intentionally general and can be loose. For small $B$, $1-2^{-B}=(\ln 2)B+O(B^2)$, so it has the required vanishing behavior, but sharper entropy--error relations are known in more specialized settings \cite{konig2009,hu2012}. The role of \cref{cor:bayes} is to convert the information-theoretic theorem into a transparent operational recovery statement, not to provide the tightest possible Bayes-risk converse.
\end{remark}

\begin{remark}[What the theorem does not say]
The theorem does not prove that an arbitrary quantum environment is close to SBS, nor that maximizing \cref{eq:objective} always finds the global optimum efficiently. It says that if a measurement tuple has empirically relevant properties---low worst-pair disagreement and high minimum entropy---and the physical state is close to an SBS reference, then the learned outputs are necessarily informative about that reference pointer variable.
\end{remark}

\section{Finite-shot learning guarantee for a finite measurement class}

A NISQ-compatible learning statement must account for finite measurement shots. The following proposition is intentionally conservative; it uses only Hoeffding concentration, a union bound, and finite-alphabet entropy continuity.

Let $\mathcal C\subset\mathfrak M^{(L)}$ be a finite candidate class of size $Q$. For each $\mathbf M\in\mathcal C$, use $n$ fresh independent preparations and record the full outcome vector $(Z_1,\ldots,Z_m)$. Let $\widehat D(\mathbf M)$ be the empirical maximum pairwise disagreement and $\widehat R(\mathbf M)$ the empirical minimum marginal entropy.

\begin{proposition}[Uniform finite-shot certificate]
\label{prop:finite}
Let $C_m=\binom{m}{2}$ and choose failure probability $\beta\in(0,1)$. Define
\begin{align}
a_n&=\sqrt{\frac{\ln(4QC_m/\beta)}{2n}},\\
t_n&=\sqrt{\frac{\ln(4QmL/\beta)}{2n}},\\
\tau_n&=\frac{L t_n}{2}.
\end{align}
Assume $\tau_n\le1-1/L$. Then with probability at least $1-\beta$, simultaneously for every $\mathbf M\in\mathcal C$,
\begin{align}
|\widehat D(\mathbf M)-D(\mathbf M)|&\le a_n,\label{eq:Duniform}\\
|\widehat R(\mathbf M)-R(\mathbf M)|&\le \cL(\tau_n).\label{eq:Runiform}
\end{align}

Therefore, if an empirical EIQL solver selects
\begin{equation}
\widehat{\mathbf M}\in\arg\max_{\mathbf M\in\mathcal C}\widehat R(\mathbf M)
\quad\text{subject to}\quad
\widehat D(\mathbf M)\le\eps-a_n,
\label{eq:empsolver}
\end{equation}
then $D(\widehat{\mathbf M})\le\eps$. Moreover, relative to the best candidate satisfying the stricter population margin $D(\mathbf M)\le\eps-2a_n$, the selected true richness is at most $2\cL(\tau_n)$ below optimal.
\end{proposition}

\begin{proof}[Proof sketch]
For each candidate and each fragment pair, disagreement is a Bernoulli random variable. Hoeffding's inequality and a union bound over $QC_m$ pair-statistics yield \cref{eq:Duniform}. For every candidate, fragment, and output symbol, Hoeffding plus a union bound over $QmL$ probabilities yields coordinate error at most $t_n$, hence marginal total variation at most $\tau_n$. The finite-alphabet continuity bound then gives \cref{eq:Runiform}. The feasibility and near-optimality statements follow by applying these two uniform bounds to the empirical optimizer and to the population comparator.
\end{proof}

This result certifies a pre-specified finite candidate set evaluated with fresh data. It does not automatically certify adaptive candidate generation in which future measurement settings depend on earlier shot outcomes; that setting requires a sequential or other adaptivity-aware generalization analysis. More broadly, this is not a sample-complexity theorem for a continuous POVM manifold. For continuous classes, covering numbers, Rademacher-type complexity, or structure-specific geometry are natural next steps \cite{banchi2023}. It does, however, make a finite-grid or calibrated-measurement EIQL experiment statistically well defined.

\section{Hardware-compatible finite-measurement protocol}

A minimal EIQL experiment can be run without a deep variational circuit:

\begin{enumerate}
\item Prepare the same open-system process repeatedly and obtain $m$ environmental fragments from each run.
\item Choose a hardware-constrained candidate local measurement tuple $\mathbf M$.
\item Measure all fragments on fresh runs and estimate $\widehat D(\mathbf M)$ and $\widehat R(\mathbf M)$.
\item Search the candidate class for maximum richness under the disagreement constraint, optionally using the lexicographic minimum-disagreement tie-break.
\item Validate the selected measurement on fresh shots and compare with a permutation null that independently shuffles fragment outcomes across physical runs.
\end{enumerate}

For single-qubit fragments, a candidate measurement can be a Bloch-axis projective measurement, implemented by one local pre-rotation and a computational-basis measurement. The quantum device supplies the state and measurement statistics; the outer search can be a pre-specified classical grid, random finite candidate set, or finite calibrated basis set while retaining \cref{prop:finite}. Adaptive procedures such as Bayesian optimization may still be useful experimentally, but require a separate sequential/generalization analysis.

Current quantum hardware has already been used to witness Quantum Darwinism and objectivity in collision models \cite{chisholm2022,saini2020}, and superconducting-circuit experiments have observed branching structures associated with Quantum Darwinism while proposing inexpensive designed local witnesses \cite{zhu2025}. EIQL's proposed distinction is not the use of quantum processors for Quantum Darwinism; it is the learning of an initially unknown local decoder from redundancy. Version 3 does not include a new hardware run; hardware validation is an optional follow-up rather than a logical premise of the framework.

\section{Numerical studies aligned with the EIQL objective}

All numerical results in this section optimize the same EIQL objective used in the theory: maximum minimum-fragment entropy subject to a worst-pair disagreement constraint. Earlier mutual-information-based exploratory simulations are not used as evidence for the theorems.

\subsection{Five-qubit collision model}

The simulated register contains one system qubit and four environment qubits,
\[
S+E_1+E_2+E_3+E_4.
\]
For each dynamical case, the hidden system basis and the \emph{common} hidden environmental basis are randomized independently of one another by $SU(2)$ rotations. The environmental basis is shared across the four fragments; it is not independently randomized for each environment qubit. The system begins in an equal superposition of its hidden pointer eigenstates and every environment qubit begins in the same hidden record state. In the stable cases, each collision uses
\begin{equation}
U_j=\exp\!\left[-i\gamma\,P_{1,S}\otimes Y_E\right],
\label{eq:collision}
\end{equation}
where $P_{1,S}$ is the hidden system pointer projector and $Y_E$ is a Pauli axis in the common hidden environmental basis. The learner is not given either basis.

The hardware class is deliberately narrow: the same one-qubit projective measurement axis is used on every fragment. Candidate axes are sampled uniformly on the Bloch sphere. This tied-decoder restriction is smaller than the full POVM class of the theorems.

The dynamical cases are stable-perfect ($\gamma=\pi/2$), stable-partial ($\gamma=\pi/3$), drift-15deg (perfect collisions plus a $15^\circ$ system rotation between collisions), scrambled-system (the system coupling projector is independently randomized at each collision), and no-collision ($\gamma=0$).

\subsection{Monte-Carlo approximation to the population frontier}

The continuous Bloch-sphere optimum is approximated by 6000 randomly sampled measurement axes per dynamical case. Therefore the numerical quantity plotted below is a Monte-Carlo lower bound on the optimum over the continuous tied-axis hardware class, not an exact population frontier.

\begin{figure}[H]
\centering
\includegraphics[width=0.88\textwidth]{eiql_v21_frontier.png}
\caption{Monte-Carlo approximation to the EIQL population agreement-richness frontier using 6000 randomly sampled Bloch axes per case. The dashed curve is the analytic independent-product binary baseline from \cref{lem:independent}. Perfect stable dynamics support one bit at very small disagreement, while the no-collision control follows the independent baseline.}
\label{fig:frontier}
\end{figure}

At $\eps=0.01$, the stable-perfect model attains one bit within the sampled class with tie-break disagreement $3.71\times10^{-4}$, while the no-collision control reaches about $0.0453$ bit. At $\eps=0.10$, stable-perfect remains at one bit, stable-partial is still infeasible, drift remains feasible and rich, and the scrambled-system model remains infeasible in the tested tied-axis class. The no-collision Monte-Carlo value is $0.2975$ bit, matching the analytic independent-product benchmark to numerical search precision.

\begin{lemma}[Independent binary baseline]
\label{lem:independent}
Suppose the fragment outputs are independent and identically distributed Bernoulli variables with $P(Z_j=1)=p$. Then
\[
D=2p(1-p),\qquad R=h_2(p).
\]
For $0\le\eps\le1/2$, the maximum richness compatible with $D\le\eps$ is
\begin{equation}
R_{\mathrm{ind}}(\eps)=
h_2\!\left(\frac{1-\sqrt{1-2\eps}}{2}\right).
\label{eq:indbaseline}
\end{equation}
\end{lemma}

\begin{proof}
For independent Bernoulli outputs, disagreement is $P(0,1)+P(1,0)=2p(1-p)$. The entropy $h_2(p)$ increases as $p$ moves from $0$ toward $1/2$. Under $2p(1-p)\le\eps$, the entropy-maximizing feasible point lies at equality and has $p=(1-\sqrt{1-2\eps})/2$, up to the entropy-symmetric replacement $p\mapsto1-p$.
\end{proof}

This lemma formalizes why agreement alone is insufficient: nearly deterministic independent outputs can have very small disagreement, but their richness is necessarily small. Genuine redundant records can lie far above this independent baseline.

\subsection{Search convergence}

Because the stable-perfect balanced model has a flat one-bit marginal-entropy maximum, the practical lexicographic tie-break is what localizes the record axis. A nested random-axis study shows systematic convergence of the best disagreement among near-maximal-richness candidates: the disagreement decreases from $4.46\times10^{-2}$ with 50 axes to $9.0\times10^{-5}$ with 20000 axes, while the angle to the hidden record axis decreases from $17.4^\circ$ to $0.77^\circ$.

\begin{figure}[H]
\centering
\includegraphics[width=0.78\textwidth]{eiql_v21_search_convergence.png}
\caption{Monte-Carlo search convergence in the stable-perfect population model. Increasing the number of randomly sampled Bloch axes improves the minimum-disagreement representative selected among near-maximal-richness candidates.}
\label{fig:convergence}
\end{figure}

\subsection{Finite-shot measurement discovery}

At $\eps=0.10$, each candidate receives 600 fresh shots, 500 candidate axes are searched, and 20 independent searches are performed for each dynamical case. Population validation is then computed on the selected axis. Table~\ref{tab:finite} reports mean $\pm$ standard deviation; the accompanying text also gives 95\% confidence-interval half-widths for the principal quantities.

\begin{table}[H]
\centering
\caption{Finite-shot EIQL at $\eps=0.10$. Values are population validation of measurements learned from finite-shot data over 20 independent searches per case.}
\label{tab:finite}
\begin{tabular}{lcccc}
\toprule
Case & Feasible & $\min_j H(Z_j)$ & $D_{\max}$ & Axis error \\
\midrule
stable-perfect & 1.00 & $1.000\pm0.000$ & $0.00194\pm0.00168$ & $3.30^\circ\pm1.41^\circ$\\
drift-15deg & 1.00 & $1.000\pm0.000$ & $0.04732\pm0.00461$ & $7.16^\circ\pm2.71^\circ$\\
no-collision & 1.00 & $0.252\pm0.022$ & $0.08089\pm0.00889$ & not physically diagnostic\\
stable-partial & 0.00 & -- & -- & --\\
scrambled-system & 0.00 & -- & -- & --\\
\bottomrule
\end{tabular}
\end{table}

For stable-perfect, the 95\% confidence-interval half-width is approximately $7.4\times10^{-4}$ for $D_{\max}$ and $0.62^\circ$ for axis error. For drift-15deg the corresponding values are about $2.0\times10^{-3}$ and $1.19^\circ$. These uncertainty estimates quantify search-to-search variability; they do not include hardware-model uncertainty because the present study is a noiseless state-vector simulation.

\subsection{Permutation null with the correct tail}

For an independent stable-perfect validation run, the learned decoder representative yielded empirical minimum entropy $0.999946$ bit and worst-pair disagreement
\[
D_{\mathrm{obs}}=0.002667
\]
on 3000 validation shots. The fragment columns were then independently permuted across physical runs 1000 times. This preserves each single-fragment marginal while destroying same-event correlations. The null disagreement had mean $0.51155$, standard deviation $0.00597$, 5th percentile $0.50233$, and median $0.51100$.

Because the EIQL signal is \emph{small} disagreement, the relevant permutation statistic is the lower tail,
\begin{equation}
p_{\mathrm{perm}}=
\frac{1+\#\{D_{\mathrm{perm}}\le D_{\mathrm{obs}}\}}{B+1}.
\label{eq:permp}
\end{equation}
No permutation among $B=1000$ produced disagreement as small as the observation, giving the corrected Monte-Carlo value
\[
p_{\mathrm{perm}}=\frac{1}{1001}\approx9.99\times10^{-4}.
\]
This is a simulation diagnostic, not a universal hypothesis test for arbitrary EIQL experiments.

\subsection{NISQ-style noise stress test}

To test whether the agreement-richness separation immediately disappears under elementary device imperfections, we repeated the stable-perfect collision experiment with a density-matrix noise model. After each system--environment collision, the participating qubit pair is subjected to a local two-qubit depolarizing channel
\begin{equation}
\rho\mapsto (1-p_2)\rho+p_2\left(\frac{I_{SE_j}}{4}\otimes\tr_{SE_j}\rho\right),
\label{eq:depol}
\end{equation}
with the tensor factor embedded at the original qubit locations. Each final fragment readout is then independently flipped with probability $q$. This is a deliberately simple NISQ-style stress model, not a calibrated model of any particular processor. For a true pairwise disagreement $D$, symmetric independent readout flips transform it to
\begin{equation}
D_{\rm obs}=2q(1-q)+(1-2q)^2D.
\label{eq:readoutdis}
\end{equation}
The population search used 5000 random tied Bloch axes at $\eps=0.10$. The tested noise grid was
\[
p_2\in\{0,0.002,0.005,0.01,0.02\},\qquad
q\in\{0,0.005,0.01,0.02,0.05\}.
\]
We then repeated representative settings over 12 independently randomized hidden-basis worlds. Each world used independent system and common-environment basis draws and 3500 candidate axes.

\begin{table}[H]
\centering
\caption{NISQ-style noise stress test over 12 independently randomized hidden-basis worlds. Richness remains one bit in this balanced binary model, so feasibility is controlled by worst-pair disagreement.}
\label{tab:noisy}
\begin{tabular}{lccccc}
\toprule
Setting & $p_2$ & $q$ & Feasible & Mean $D_{\max}$ & Axis error \\
\midrule
clean & 0 & 0 & 1.00 & $0.00013\pm0.00013$ & $0.82^\circ\pm0.44^\circ$\\
 mild & 0.005 & 0.01 & 1.00 & $0.02951\pm0.00019$ & $0.96^\circ\pm0.59^\circ$\\
 moderate & 0.01 & 0.02 & 1.00 & $0.05764\pm0.00029$ & $1.28^\circ\pm0.72^\circ$\\
 readout-heavy & 0.002 & 0.05 & 1.00 & $0.09834\pm0.00009$ & $0.90^\circ\pm0.36^\circ$\\
 beyond-$\eps$ & 0.02 & 0.05 & 0.00 & $0.12676\pm0.00023$ & --\\
\bottomrule
\end{tabular}
\end{table}

\begin{figure}[H]
\centering
\includegraphics[width=0.78\textwidth]{eiql_noisy_nisq_disagreement.png}
\caption{Worst-pair disagreement under local two-qubit depolarization and symmetric readout flips. The dashed line is the EIQL feasibility tolerance $\eps=0.10$. The stress test shows a finite operating region rather than noise immunity: sufficiently large combined noise pushes the stable record beyond the chosen tolerance.}
\label{fig:noisy}
\end{figure}

The clean-to-moderate regime remains well separated from the $\eps=0.10$ boundary, while $5\%$ readout flips become marginal and fail once combined with stronger collision depolarization. Across feasible settings, the learned axis remains within roughly $1$--$1.3^\circ$ on average of the hidden record axis. These results support only a limited claim: the EIQL objective does not collapse under this simple local noise model. They do not substitute for a hardware experiment, calibrated crosstalk/readout models, time-correlated drift, or a scaling study.

Taken together, the noiseless and noisy collision simulations show that the EIQL objective is operational and behaves as designed on a small model. They do not establish computational advantage, favorable asymptotic scaling, or device-independent robustness.

\section{External-architecture benchmark: Chen et al. photonic Quantum Darwinism}

To test EIQL on a structure not invented for this manuscript, we constructed an unknown-decoder benchmark from the six-photon Quantum-Darwinism experiment of Chen et al. \cite{chen2019}. Their system consists of one photon representing the system and five photons representing environmental records. Two published interaction settings are
\[
\boldsymbol\theta_A=(180^\circ,180^\circ,180^\circ,180^\circ,180^\circ)
\]
and
\[
\boldsymbol\theta_B=(180^\circ,180^\circ,180^\circ,72^\circ,100^\circ),
\]
so setting $A$ contains five ideal binary records while setting $B$ contains three ideal records and two intrinsically weaker records.

The EIQL modification is deliberately blind: each environmental fragment is independently conjugated by a hidden local $SU(2)$ rotation, and the learner is not given the system label, pointer value, interaction angle, or hidden readout basis. Separate local projective decoders are learned from environmental redundancy alone. The hidden rotations and pointer labels are revealed only after training for evaluation. This converts a published Quantum-Darwinism record architecture into a measurement-learning task without claiming that the original experiment itself performed EIQL.

A multi-start population search was repeated over 12 independently hidden-basis worlds. Table~\ref{tab:chen} compares the learned worst-pair disagreement with the physical floor obtained from the oracle Helstrom record quality.

\begin{table}[H]
\centering
\caption{EIQL on the Chen et al. record architecture. The fidelity-proxy rows use an isotropic white-noise approximation matched to the reported full-state fidelities; they are not a reconstruction of the laboratory noise channel.}
\label{tab:chen}
\begin{tabular}{lccc}
\toprule
Setting & Mean $\min_j H(Z_j)$ & Learned $D_{\max}$ & Oracle floor \\
\midrule
$A$ ideal & $1.00000$ & $0.00193\pm0.00093$ & $0$ \\
$B$ ideal & $0.99910$ & $0.27689\pm0.00496$ & $0.27487$ \\
$A$ fidelity proxy & $1.00000$ & $0.07987\pm0.01624$ & $0.07162$ \\
$B$ fidelity proxy & $0.99915$ & $0.35624\pm0.04524$ & $0.34279$ \\
\bottomrule
\end{tabular}
\end{table}

\begin{figure}[H]
\centering
\includegraphics[width=0.78\textwidth]{eiql_chen_multistart_vs_oracle.png}
\caption{Learned worst-pair disagreement on the Chen et al. architecture compared with the oracle physical floor. In the ideal strong-record setting the residual is search error; in the mixed-quality setting the dominant disagreement is intrinsic to the two weak records.}
\label{fig:chenoracle}
\end{figure}

The mixed setting also provides a fragment-quality test. EIQL was allowed to rank every three-fragment subset. The correct high-quality subset---photons $2,3,4$, corresponding to the three $180^\circ$ records---was selected in $12/12$ population worlds and $10/10$ finite-shot worlds at 700 shots per estimated statistic. The learner was not given which fragments were strong. This result is stronger than merely recovering one hidden basis: the same redundancy criterion identifies which environmental records are trustworthy.

A finite-shot version gives mean validated disagreement $0.01791$ for $A$ and $0.28420$ for $B$ at 700 shots per estimated statistic. The latter remains close to the physical $B$-setting floor. These experiments are simulations based on the published architecture, not reanalysis of raw photonic events.

\section{Environment-only pair-moment implementation and resource benchmark}

The preceding searches treat EIQL as optimization over candidate decoders. In the binary conditionally independent qubit-record model, the redundancy structure also yields a direct spectral implementation. Let $\mathbf r_{jx}$ be the Bloch vector of fragment $j$ conditioned on the balanced hidden branch $X=x\in\{0,1\}$. Measuring Pauli components $a,b\in\{x,y,z\}$ on fragments $i,j$ gives a centered cross-correlation matrix
\begin{equation}
C_{ij}=\frac14\bigl(\mathbf r_{i0}-\mathbf r_{i1}\bigr)
                 \bigl(\mathbf r_{j0}-\mathbf r_{j1}\bigr)^T.
\label{eq:pairmoment}
\end{equation}
Thus $C_{ij}$ is rank one in the ideal model, and its leading singular directions recover the local binary record directions up to the free sign/output relabeling. For equal priors those directions coincide with the local Helstrom discrimination axes. We use this identity as a task-specific implementation device, not as a novelty claim; related covariance and spectral techniques should be checked before attributing originality to the factorization itself.

To estimate all pair moments without full state tomography, we use a Pauli orthogonal-array schedule in which every pair of fragments encounters all nine ordered Pauli-basis pairs. Five environment fragments require 27 distinct settings. For comparison, a stronger task-specific baseline with direct access to $S$ reconstructs every $S$--$E_j$ $3\times3$ correlation matrix in 9 settings, whereas full Pauli tomography of the six-qubit $S E_1\cdots E_5$ state requires $3^6=729$ settings. The comparison is intentionally not ``EIQL versus only full tomography'': the 9-setting system-assisted baseline is included to avoid an inflated resource claim.

At equal total copy budget, the environment-only spectral EIQL estimator and the system-assisted baseline were evaluated over 40 independently hidden-basis worlds. At 6912 copies per method, the results are:

\begin{table}[H]
\centering
\caption{Equal-copy comparison between environment-only EIQL and a stronger system-assisted correlation estimator. Axis error is measured against the hidden local Helstrom directions after training.}
\label{tab:resource}
\begin{tabular}{lcccc}
\toprule
Setting & EIQL axis error & $S$-assisted axis error & EIQL $D_{\max}$ & $S$-assisted $D_{\max}$ \\
\midrule
Chen $A$ strong & $1.63^\circ$ & $2.32^\circ$ & $0.00092$ & $0.00194$ \\
Chen $B$ mixed & $2.08^\circ$ & $2.65^\circ$ & $0.27536$ & $0.27592$ \\
\bottomrule
\end{tabular}
\end{table}

\begin{figure}[H]
\centering
\includegraphics[width=0.82\textwidth]{eiql_vs_tomography_equal_budget.png}
\caption{Decoder-axis error versus equal total copy budget. The environment-only estimator aggregates redundancy across all fragment pairs, while the comparison baseline has direct access to the system. The numerical advantage in this model is not claimed to be universal.}
\label{fig:equalbudget}
\end{figure}

Readout-flip stress at 256 shots per EIQL setting remains close to the corresponding noise-imposed physical floor: in the strong-record case at $q=0.05$, the mean axis error is $2.06^\circ$ and the observed disagreement is $0.09612$ versus an oracle observed floor of $0.095$. An independent-fragment null at the same shot count gives mean spectral pair signal $0.09082$, compared with $0.99931$ for same-event redundant records, while balanced independent fragments have population disagreement $0.5$.

The measurement-setting comparison has very different scaling from full Pauli tomography. For the orthogonal-array construction used here, 9 settings support up to four environment fragments and 27 settings support up to thirteen, while full Pauli tomography of the system plus $m$ fragments scales as $3^{m+1}$.

\begin{figure}[H]
\centering
\includegraphics[width=0.78\textwidth]{eiql_vs_tomography_settings_scaling.png}
\caption{Distinct measurement settings for the task-specific environment-only pair-moment design, the system-assisted pair-correlation baseline, and full Pauli tomography. This figure compares settings, not total runtime or a universal minimax sample complexity.}
\label{fig:settings}
\end{figure}

The defensible resource conclusion is narrow: in this structured binary-record task, EIQL can recover hidden environmental decoders without direct system access and without reconstructing the full state. Full tomography is substantially more measurement-setting intensive, but a task-specific method with system access can use fewer settings than EIQL. Therefore Version 3 does not claim universal tomography superiority; it identifies the operational tradeoff created by removing system access and exploiting inter-fragment redundancy.

\section{Novelty positioning and neighboring literature}

The publication claim must be narrower than the v1 wording. The following distinctions are central.

\begin{table}[H]
\centering
\caption{What is established prior art, and what EIQL proposes to add.}
\begin{tabularx}{\textwidth}{>{\raggedright\arraybackslash}p{0.29\textwidth}X}
\toprule
Established direction & Relation to EIQL \\
\midrule
Environment as witness / Quantum Darwinism \cite{ollivier2004,ollivier2005} & Establishes redundant environmental records and the selection of pointer information. EIQL does not claim this mechanism as new.\\
SBS and strong Quantum Darwinism \cite{horodecki2015,le2019} & Supplies the structural model used in the exact and robust theorems.\\
Uniqueness of redundant imprints \cite{fu2021} & Shows that pointer-observable uniqueness under redundant imprint conditions is already a dedicated theorem topic. EIQL therefore does not market Theorem~\ref{thm:exact} as the primary novelty.\\
Observer consensus \cite{touil2025} & Proves agreement among observers with informative fragments. EIQL asks the different operational question of learning the unknown local measurement/decoder that generates a rich agreeing variable.\\
Observer resource dependence \cite{feller2023} & Shows that accessible fragments and measurement resources matter for reconstructing the classical picture, reinforcing the need for an explicit hardware-constrained learning formulation.\\
Experimental Quantum Darwinism \cite{chen2019,saini2020,chisholm2022,zhu2025} & Establishes photonic and processor-based generation/witnessing of redundant records, and in Zhu et al. a designed inexpensive local witness. EIQL uses these works as physical motivation but treats the environmental decoder as initially unknown and learned.\\
Petz recovery of Darwinian information \cite{torvinen2026} & Studies a specified recovery map from fragments. EIQL instead optimizes an unknown decoder from repeated data under an explicit self-supervised objective.\\
Gacs--Korner common information \cite{gacs1973,csirmaz2023} & Provides the classical mathematical archetype: the richest deterministic common variable. EIQL embeds a related idea into quantum measurement selection and hardware constraints.\\
Quantum self-supervised learning \cite{jaderberg2021} & Establishes that self-supervised objectives already exist in quantum learning. EIQL's distinction is the physical unknown-decoder task induced by environmental redundancy, not self-supervision by itself.\\
Learning unknown quantum measurements \cite{cheng2016} & Learns an externally specified unknown target measurement from training information and studies sample complexity of that target class. EIQL has no target measurement labels; redundancy across same-event fragments generates the training criterion.\\
Unsupervised observable discovery \cite{sadoune2023} & Learns interpretable characteristic observables/order parameters from quantum measurement data without Hamiltonian labels. EIQL instead searches for local decoders whose outputs are redundantly aligned across independently accessible environmental fragments.\\
Quantum-learning statistical complexity \cite{banchi2023} & Highlights data, copy, and model complexity in quantum learning. EIQL's finite-class proposition is a deliberately simple copy-complexity certificate, while adaptive and continuous decoder classes remain open.\\
Redundancy versus consensus \cite{chisholm2023} & Shows that these notions need not coincide. EIQL therefore defines its operational objective directly through aligned worst-pair decoder disagreement and minimum fragment richness rather than treating foundations terminology as interchangeable.\\
\bottomrule
\end{tabularx}
\end{table}

Accordingly, the strongest defensible novelty hypothesis is:

\begin{quote}
EIQL is a quantum measurement-learning framework that extends the principle of multiview agreement to quantum environmental records, where the representation must itself be physically instantiated through a learned measurement. Its specific learning signal is same-event redundancy across independently accessible fragments, and near SBS states the resulting agreement-richness objective carries quantitative pointer-information guarantees.
\end{quote}

The novelty claim is intentionally at the level of a quantum learning framework/problem formulation, not a new universal paradigm. Direct literature comparison found strong neighboring results in Quantum Darwinism, consensus, unknown-measurement learning, quantum self-supervision, and unsupervised observable discovery, but not the same combination of unknown local decoders, environment-only same-event redundancy as supervision, and pointer-information guarantees. Because terminology differs across fields, this remains a claim for external referee scrutiny rather than an assertion of exhaustive uniqueness.

\section{What is quantum in EIQL?}

The redundantly objective variable in an SBS state is classical. The quantum content lies in three places:

\begin{enumerate}
\item \textbf{The records are quantum states.} Before measurement, each fragment is described by a density operator and incompatible observables are not simultaneously available as a classical feature vector.
\item \textbf{The representation is a POVM.} Learning means choosing a physically implementable measurement/decoder, not merely selecting a coordinate from an already classical data table.
\item \textbf{The redundancy is generated by quantum dynamics.} Entangling system-environment interactions, decoherence, and einselection determine which record can proliferate.
\end{enumerate}

Once measurement outcomes have been produced, the disagreement and entropy statistics are classical. This is why Version 3 uses the narrower subtitle ``Self-Supervised Measurement Discovery from Redundant Quantum Records'' rather than language suggesting that an arbitrary quantum state itself becomes an objective classical fact.

\section{Relation to variational QML and barren plateaus}

EIQL is not defined by a parameterized circuit. In the finite candidate protocol of \cref{prop:finite}, learning is discrete measurement selection, and the statistical difficulty is shot complexity rather than gradient concentration. In the single-qubit simulation, each candidate is just a local projective basis.

This is a design-level avoidance of one common source of barren plateaus, not a universal trainability guarantee. A future EIQL implementation could itself parameterize deep measurement circuits, in which case barren plateaus or other nonconvex pathologies may return. The correct claim is therefore:

\begin{quote}
Barren plateaus are not intrinsic to the minimal EIQL formulation because EIQL does not require deep variational quantum-circuit training.
\end{quote}

The broader barren-plateau literature remains directly relevant if richer decoder classes are implemented variationally \cite{mcclean2018,larocca2024}.

\section{Interpretability and downstream use}

EIQL produces an operationally interpretable object: a local decoder behavior (defined modulo equivalence on the relevant record states) and a common classical outcome variable. For each learned representation one can report:

\begin{itemize}
\item the decoder behavior implemented on each fragment, reported modulo operational equivalence on the relevant record states;
\item the worst-pair disagreement $D_\rho$;
\item the minimum fragment entropy $R_\rho$;
\item the full frontier $\mathcal O_{\eps,L}^{\mathfrak M}$ as a reliability--richness tradeoff;
\item ablations showing which fragments are necessary for the common record.
\end{itemize}

This form of interpretability is physical and operational rather than post-hoc. It does not imply causality: a redundantly shared variable can still be a nuisance variable if the physical process broadcasts that nuisance. Any downstream application therefore needs independent task validation.

A classical multimodal analogue---for example blood, MRI, and ECG measurements from one patient---can use the same redundancy principle, but such a model is not quantum merely because it was inspired by EIQL. The genuinely quantum version requires quantum fragments and learned POVMs.

\section{Limitations and falsification criteria}

EIQL should be considered falsified or substantially weakened as a useful framework if one or more of the following persist under stronger analysis:

\begin{enumerate}
\item \textbf{No novelty after direct comparison.} Existing work may already contain the same hardware-constrained, unknown-decoder learning problem with comparable robust and finite-sample guarantees.
\item \textbf{Objective degeneracy.} Large classes of non-Darwinian states may achieve the same agreement-richness frontier as SBS-like states without meaningful record structure.
\item \textbf{Sample and setting complexity beyond the structured qubit model.} The pair-moment benchmark in \cref{fig:settings} is favorable relative to full Pauli tomography but exploits binary, conditionally independent qubit records. General POVM classes or higher-dimensional fragments may require substantially more copies and settings.
\item \textbf{No real-hardware EIQL result.} The simple stress test in \cref{fig:noisy} and the Chen/resource benchmarks are simulations. Crosstalk, coherent calibration errors, time-correlated drift, leakage, and fragment non-independence may create spurious agreement or erase the operational gap between stable and control dynamics. A hardware run is useful validation but is outside the scope frozen for Version 3.
\item \textbf{Theorem vacuity.} If realistic states cannot be related to an SBS reference with useful $\eta$, the robust theorem may remain mathematically correct but physically uninformative.
\end{enumerate}

The no-collision control in \cref{fig:frontier} already illustrates why both low disagreement and high richness are necessary. A nearly deterministic measurement can manufacture agreement from independent product fragments, but it cannot simultaneously produce one bit of common richness at small $\eps$.

\section{Open theoretical problems}

The most important next problems are now well defined.

\begin{enumerate}
\item \textbf{Continuous measurement classes.} Replace the finite-class bound in \cref{prop:finite} by covering-number or geometry-aware bounds for POVM manifolds.
\item \textbf{Approximate structure without an SBS oracle.} Reverse part of the present logic: derive approximate latent broadcast structure from operational EIQL properties of $\rho$ itself, such as multiple rich agreeing local decoders or approximate conditional-independence/recoverability conditions. This would turn EIQL from a learning layer conditioned on nearby SBS structure into a possible diagnostic of Darwinian structure.
\item \textbf{General independent null ceilings.} Extend \cref{lem:independent} from i.i.d. Bernoulli fragments to non-identically distributed and $L$-ary independent fragments, yielding a broader operational null ceiling for agreement-richness witnesses.
\item \textbf{Fragment correlations.} Quantify how violations of strong conditional independence modify \cref{thm:robust}. Strong Quantum Darwinism and SBS literature shows that this distinction is nontrivial \cite{le2019}.
\item \textbf{Adaptive and collective measurements.} Determine when local independent POVMs are statistically suboptimal relative to adaptive or collective fragment strategies.
\item \textbf{General resource theory.} Extend the structured pair-moment benchmark into copy-, setting-, and runtime-complexity bounds for general fragment dimensions, decoder classes, and access models (environment-only versus direct system access).
\item \textbf{Operational equivalence and identifiability at nonzero $\eps$.} Richness can be flat over many physically distinct POVMs. Characterize when the minimum-disagreement tie-break converges to a unique decoder equivalence class on the relevant record family.
\item \textbf{Downstream learning.} Determine when an EIQL-discovered variable improves prediction, control, or sensing relative to tomography-first and supervised alternatives.
\end{enumerate}

\section{Proposed definition of EIQL}

A concise formal definition of the research program is:

\begin{definition}[Environment-Induced Quantum Learning]
An EIQL problem consists of (i) repeated physical access to multiple quantum fragments generated by a common underlying process, (ii) a hardware-constrained class of local measurements or instruments, and (iii) an unlabeled objective that rewards information which is simultaneously rich and redundantly recoverable across the fragments. A learned representation is a physical decoder/measurement, and its statistical quality is evaluated operationally from repeated measurement outcomes.
\end{definition}

For the present theory, the canonical objective is \cref{eq:objective}. The exact SBS theorem identifies its zero-error structure. The robust theorem states when high empirical-quality outputs imply pointer recoverability. The finite-class proposition turns the objective into a finite-shot learning problem.

\section{Conclusion}

Version 3 frames EIQL as a specific quantum measurement-learning framework rather than a new universal learning paradigm. Quantum Darwinism, SBS, redundant-imprint uniqueness, observer consensus, and experimental witnessing of Darwinian records are established prior art. Classical common-information and multiview-agreement principles also make clear that ``learn what multiple views agree on'' is not itself the contribution.

The proposed learning problem begins where those ingredients leave an operational gap: the observer is not handed the environmental decoder. It must choose a physically implementable local measurement whose outcomes are both rich and redundant across independently accessible fragments generated by the same event. Under exact SBS, exact agreement identifies deterministic coarse-grainings of the pointer variable; near SBS, the robust theorem bounds residual pointer uncertainty; and finite pre-specified decoder classes admit explicit finite-shot guarantees. These statements concern pointer information and decoder behavior on the relevant record family, not uniqueness of a POVM operator on the full Hilbert space.

The numerical evidence now has three layers. First, theory-matched collision simulations reproduce the analytic independent-product null ceiling, separate stable redundant dynamics from controls, pass a lower-tail permutation test, and retain a finite operating region under a simple depolarizing-plus-readout noise model. Second, a benchmark based on the published Chen et al. photonic architecture learns independently hidden local decoders and identifies the strong-record subset without system labels. Third, a pair-moment spectral implementation recovers decoder directions without direct system access and is compared against both a stronger system-assisted task-specific baseline and full tomography in measurement-setting count. The comparison is deliberately qualified: EIQL avoids full-state reconstruction, but it is not universally cheaper than every task-specific method that is allowed to interrogate the system directly.

The first-paper scope therefore stops at theory plus reproducible simulation. A real-hardware EIQL run would be useful follow-up evidence, but it is not required for the logical claims made here. The remaining scientific questions are whether the formulation survives broader literature scrutiny, how far the guarantees extend beyond nearby SBS/product-record structure, and what resource advantages or disadvantages appear under realistic access constraints.

\appendix

\section{Additional proof details}

\subsection{Monotonicity of $\cL$}
For $L\ge2$,
\[
\cL(t)=h_2(t)+t\log_2(L-1)
\]
has derivative
\[
\cL'(t)=\log_2\frac{(1-t)(L-1)}{t},
\]
so it is nondecreasing on $0\le t\le1-1/L$. This is the range used in \cref{thm:robust} and \cref{prop:finite}.

\subsection{Why the global error-indicator proof is needed}
The v1 proof bounded $H(Z_j\mid X=x)$ by a function of the conditional disagreement $\delta_x$ and then averaged. That step is unsafe when some $\delta_x>1-1/L$, even if the average disagreement is small. Version 2 instead defines the most likely output $g_j(x)$, proves the global classification error $e\le\delta$, and applies the single global error indicator $E$. No pointwise small-error condition is required.

\section{Reproducibility details for Version 3}

The primary theory-matched collision simulation is implemented in \texttt{eiql\_v21\_simulation.py}. The five-qubit register contains one system qubit and four environment qubits. The system and common environmental bases are independently Haar-randomized by $SU(2)$ draws; the environmental basis is shared across all four environment qubits. Population frontiers use 6000 uniformly random Bloch axes per dynamical case. The search-convergence study uses a nested set of up to 20000 axes.

For finite-shot learning at $\eps=0.10$, each candidate receives 600 fresh shots; 500 candidate axes are searched; 20 independent searches are performed per dynamical case. The lexicographic near-optimality tolerance is $0.03$ bit. The permutation null uses 3000 validation shots and $B=1000$ independent column permutations, with significance computed from the lower tail as in \cref{eq:permp}.

Accompanying result files are:
\begin{itemize}
\item \path{eiql_v21_frontier_mc.csv} and \path{eiql_v21_independent_baseline.csv};
\item \path{eiql_v21_finite_summary.csv} and \path{eiql_v21_permutation_summary.csv};
\item \path{eiql_v21_search_convergence.csv};
\item Noisy-NISQ outputs: \path{eiql_noisy_nisq_results.csv};\newline \path{eiql_noisy_nisq_multiseed_summary.csv};\newline \path{eiql_noisy_nisq_simulation.py}.
\item Chen-architecture benchmark: \path{eiql_chen2019_benchmark.py};\newline \path{eiql_chen_multistart_summary.csv};\newline \path{eiql_chen_finite_stat_summary.csv}.
\item Resource benchmark: \path{eiql_vs_tomography_resource_benchmark.py};\newline \path{eiql_vs_tomography_summary.csv};\newline \path{eiql_vs_tomography_noise_summary.csv};\newline \path{eiql_vs_tomography_scaling.csv}.
\end{itemize}

\section{Version 3 change summary}

\begin{enumerate}
\item Reframed EIQL consistently as a \emph{quantum measurement-learning framework/problem formulation}, not a new universal learning paradigm.
\item Added explicit discussion of classical common-information/multiview relatives and existing quantum self-supervised learning; the quantum distinction is that the representation is physically created by a learned POVM.
\item Preserved the Version 2.1 theorem repairs: pointer-information rather than POVM identifiability, free output relabelings, alphabet qualification, robust global-error proof, finite-class finite-shot guarantee, and the independent-product null.
\item Added the external-architecture benchmark based on Chen et al. (2019), including independently hidden local record bases, mixed record quality, subset discovery, and finite-shot validation.
\item Added the environment-only pair-moment/spectral implementation and an equal-copy comparison with a stronger system-assisted correlation baseline.
\item Added measurement-setting scaling against full Pauli tomography while explicitly avoiding a universal tomography-superiority claim.
\item Froze the first-paper scope at theory plus simulation. A real-hardware run is documented as future validation rather than a prerequisite for Version 3.
\item Updated novelty positioning against Chen et al., Saini--Behera, Zhu et al., Cheng et al., Sadoune et al., Jaderberg et al., and the established Quantum-Darwinism/SBS literature.
\end{enumerate}

\begin{thebibliography}{99}

\bibitem{zurek2003}
W. H. Zurek, ``Decoherence, einselection, and the quantum origins of the classical,'' \emph{Rev. Mod. Phys.} 75, 715 (2003). \href{https://arxiv.org/abs/quant-ph/0105127}{arXiv:quant-ph/0105127}.

\bibitem{ollivier2004}
H. Ollivier, D. Poulin, and W. H. Zurek, ``Objective Properties from Subjective Quantum States: Environment as a Witness,'' \emph{Phys. Rev. Lett.} 93, 220401 (2004). \href{https://arxiv.org/abs/quant-ph/0307229}{arXiv:quant-ph/0307229}.

\bibitem{ollivier2005}
H. Ollivier, D. Poulin, and W. H. Zurek, ``Environment as a Witness: Selective Proliferation of Information and Emergence of Objectivity in a Quantum Universe,'' \emph{Phys. Rev. A} 72, 042113 (2005). \href{https://arxiv.org/abs/quant-ph/0408125}{arXiv:quant-ph/0408125}.

\bibitem{horodecki2015}
R. Horodecki, J. K. Korbicz, and P. Horodecki, ``Quantum Origins of Objectivity,'' \emph{Phys. Rev. A} 91, 032122 (2015). \href{https://arxiv.org/abs/1312.6588}{arXiv:1312.6588}.

\bibitem{le2019}
T. P. Le and A. Olaya-Castro, ``Strong Quantum Darwinism and Strong Independence are Equivalent to Spectrum Broadcast Structure,'' \emph{Phys. Rev. Lett.} 122, 010403 (2019). \href{https://arxiv.org/abs/1803.08936}{arXiv:1803.08936}.

\bibitem{korbicz2021}
J. K. Korbicz, ``Roads to objectivity: Quantum Darwinism, Spectrum Broadcast Structures, and strong quantum Darwinism -- a review,'' \emph{Quantum} 5, 571 (2021). \href{https://arxiv.org/abs/2007.04276}{arXiv:2007.04276}.

\bibitem{fu2021}
H.-F. Fu, ``Uniqueness of the Observable Leaving Redundant Imprints in the Environment in the Context of Quantum Darwinism,'' \emph{Phys. Rev. A} 103, 042210 (2021). \href{https://arxiv.org/abs/2010.14131}{arXiv:2010.14131}.

\bibitem{touil2025}
A. Touil, B. Yan, and W. H. Zurek, ``Consensus About Classical Reality in a Quantum Universe,'' (2025). \href{https://arxiv.org/abs/2503.14791}{arXiv:2503.14791}.

\bibitem{feller2023}
A. Feller, B. Roussel, A. Pontlevy, and P. Degiovanni, ``Testing quantum Darwinism dependence on observers' resources,'' (2023). \href{https://arxiv.org/abs/2306.14745}{arXiv:2306.14745}.

\bibitem{chisholm2022}
D. A. Chisholm, G. Garc\'ia-P\'erez, M. A. C. Rossi, S. Maniscalco, and G. M. Palma, ``Witnessing Objectivity on a Quantum Computer,'' \emph{Quantum Sci. Technol.} 7, 015022 (2022). \href{https://arxiv.org/abs/2110.06243}{arXiv:2110.06243}.

\bibitem{zhu2025}
Z. Zhu et al., ``Observation of Quantum Darwinism and the Origin of Classicality with Superconducting Circuits,'' (2025). \href{https://arxiv.org/abs/2504.00781}{arXiv:2504.00781}.

\bibitem{torvinen2026}
J. Torvinen, E. Keski-Vakkuri, and N. Pranzini, ``Quantum Darwinism and the quality of Petz recovery,'' (2026). \href{https://arxiv.org/abs/2605.06848}{arXiv:2605.06848}.

\bibitem{gacs1973}
P. G\'acs and J. K\"orner, ``Common information is far less than mutual information,'' \emph{Problems of Control and Information Theory} 2(2), 149--162 (1973).

\bibitem{csirmaz2023}
L. Csirmaz, ``A short proof of the Gacs--Korner theorem,'' (2023). \href{https://arxiv.org/abs/2306.14718}{arXiv:2306.14718}.

\bibitem{winter2016}
A. Winter, ``Tight uniform continuity bounds for quantum entropies: conditional entropy, relative entropy distance and energy constraints,'' \emph{Commun. Math. Phys.} 347, 291--313 (2016). \href{https://arxiv.org/abs/1507.07775}{arXiv:1507.07775}.

\bibitem{chen2019}
M.-C. Chen, H.-S. Zhong, Y. Li, D. Wu, X.-L. Wang, L. Li, N.-L. Liu, C.-Y. Lu, and J.-W. Pan, ``Emergence of classical objectivity of quantum Darwinism in a photonic quantum simulator,'' \emph{Science Bulletin} 64, 580--585 (2019). \href{https://doi.org/10.1016/j.scib.2019.03.032}{doi:10.1016/j.scib.2019.03.032}; \href{https://arxiv.org/abs/1808.07388}{arXiv:1808.07388}.

\bibitem{saini2020}
R. Saini and B. K. Behera, ``Experimental Realization of Quantum Darwinism State on Quantum Computers,'' (2020). \href{https://arxiv.org/abs/2012.07562}{arXiv:2012.07562}.

\bibitem{jaderberg2021}
B. Jaderberg, L. W. Anderson, W. Xie, S. Albanie, M. Kiffner, and D. Jaksch, ``Quantum Self-Supervised Learning,'' \emph{Quantum Sci. Technol.} 7, 035005 (2022). \href{https://arxiv.org/abs/2103.14653}{arXiv:2103.14653}.

\bibitem{ciliberto2018}
C. Ciliberto, M. Herbster, A. D. Ialongo, M. Pontil, A. Rocchetto, S. Severini, and L. Wossnig, ``Quantum machine learning: a classical perspective,'' \emph{Proc. R. Soc. A} 474, 20170551 (2018). \href{https://arxiv.org/abs/1707.08561}{arXiv:1707.08561}.

\bibitem{cheng2016}
H.-C. Cheng, M.-H. Hsieh, and P.-C. Yeh, ``The Learnability of Unknown Quantum Measurements,'' \emph{Quantum Information and Computation} 16, 615--656 (2016). \href{https://arxiv.org/abs/1501.00559}{arXiv:1501.00559}.

\bibitem{sadoune2023}
N. Sadoune, G. Giudici, K. Liu, and L. Pollet, ``Unsupervised Interpretable Learning of Phases From Many-Qubit Systems,'' \emph{Phys. Rev. Research} 5, 013082 (2023). \href{https://arxiv.org/abs/2208.08850}{arXiv:2208.08850}.

\bibitem{banchi2023}
L. Banchi, J. L. Pereira, S. T. Jose, and O. Simeone, ``Statistical Complexity of Quantum Learning,'' (2023; journal version 2025). \href{https://arxiv.org/abs/2309.11617}{arXiv:2309.11617}.

\bibitem{chisholm2023}
D. A. Chisholm, L. Innocenti, and G. M. Palma, ``The meaning of redundancy and consensus in quantum objectivity,'' \emph{Quantum} 7, 1074 (2023). \href{https://arxiv.org/abs/2211.09150}{arXiv:2211.09150}.

\bibitem{mcclean2018}
J. R. McClean, S. Boixo, V. N. Smelyanskiy, R. Babbush, and H. Neven, ``Barren plateaus in quantum neural network training landscapes,'' \emph{Nat. Commun.} 9, 4812 (2018). \href{https://arxiv.org/abs/1803.11173}{arXiv:1803.11173}.

\bibitem{larocca2024}
M. Larocca et al., ``A Review of Barren Plateaus in Variational Quantum Computing,'' (2024; revised 2025). \href{https://arxiv.org/abs/2405.00781}{arXiv:2405.00781}.

\bibitem{konig2009}
R. K\"onig, R. Renner, and C. Schaffner, ``The Operational Meaning of Min- and Max-Entropy,'' \emph{IEEE Trans. Inf. Theory} 55, 4337--4347 (2009). \href{https://arxiv.org/abs/0807.1338}{arXiv:0807.1338}.

\bibitem{hu2012}
B.-G. Hu and H.-J. Xing, ``Analytical Bounds between Entropy and Error Probability in Binary Classifications,'' (2012). \href{https://arxiv.org/abs/1205.6602}{arXiv:1205.6602}.

\end{thebibliography}


\end{document}
